% ============================================================================
% KAPITEL 4: Diskrete Informations-Lagrange-Funktional
% ============================================================================

\chapter{Diskretes Informations-Lagrange-Funktional}
\label{ch:lagrange_funktional}

\section{Einleitung}
\label{sec:lagrange_einleitung}

Das diskrete Informations-Lagrange-Funktional ist die mathematische Grundlage der IWT. Alle Größen sind diskrete Sequenzen mit Zeitindex \( n \) und Raumindex \( k \). Die Variation erfolgt über diskrete Ableitungen, nicht über Differentiale.

Im informationsbasierten Rahmen ersetzt dieses Funktional:

\begin{itemize}
    \item die Newtonsche Bewegungsgleichung durch rekursive Update-Regeln,
    \item die Maxwell-Gleichungen durch diskrete Weber-Kopplungen,
    \item die Schrödinger-Gleichung durch diskrete Bohm-Struktur,
    \item die geometrische Gravitation durch emergente diskrete Metrik.
\end{itemize}

Es bildet die mathematische Grundlage der gesamten Theorie und zeigt, dass kontinuierliche Gleichungen als Grenzfälle einer tieferen diskreten Informationsdynamik erscheinen.

\section{Grundidee des diskreten Variationsprinzips}
\label{sec:varia_prinzip}

Die diskrete IWT geht von folgenden Prinzipien aus:

\begin{enumerate}
    \item Der physikalische Zustand ist eine diskrete Informationsverteilung \( I_k^{(n)} \).
    \item Information ist eine diskret erhaltene Größe.
    \item Dynamik ist rekursive Umlagerung von Information.
    \item Lokale und globale Dynamik sind komplementär.
\end{enumerate}

Daraus folgt, dass die Dynamik durch ein diskretes Funktional beschrieben werden muss, das sowohl lokale als auch globale Informationsstrukturen berücksichtigt.

\section{Das diskrete Informations-Lagrange-Funktional}
\label{sec:lagrange_funktional_def}

Das diskrete Informations-Lagrange-Funktional lautet allgemein:

\[
L_I \bigl[\{I_k^{(n)}\}\bigr] = \sum_{k=1}^{N} \mathcal{F}_k \Bigl( I_k^{(n)}, \Delta I_k^{(n)}, \delta I_k^{(n)} \Bigr) \Delta V_k.
\]

Hierbei ist:

\begin{itemize}
    \item \( k = 1, \dots, N \) – Index der diskreten Informationszelle (Knoten),
    \item \( I_k^{(n)} \in \mathbb{R}^+ \) – Informationswert am Knoten \( k \) zum Zeitindex \( n \),
    \item \( \Delta I_k^{(n)} \) – diskreter räumlicher Gradient (Nachbarschaftsdifferenzen),
    \item \( \delta I_k^{(n)} = I_k^{(n)} - I_k^{(n-1)} \) – diskrete zeitliche Änderung,
    \item \( \Delta V_k \) – diskretes Volumenelement (knotenspezifisch),
    \item \( \mathcal{F}_k \) – diskrete Lagrangedichte am Knoten \( k \).
\end{itemize}

\subsection{Diskrete räumliche Gradienten}
\label{sub:diskrete_gradienten}

Der diskrete Gradient am Knoten \( k \) ist definiert als:

\[
\Delta I_k^{(n)} = \sum_{l \in \mathcal{N}(k)} w_{kl} \bigl( I_l^{(n)} - I_k^{(n)} \bigr),
\]

mit:

\begin{itemize}
    \item \( \mathcal{N}(k) \) – Nachbarknoten von \( k \),
    \item \( w_{kl} \) – Kopplungsgewichte (normiert: \( \sum_l w_{kl} = 1 \)).
\end{itemize}

\subsection{Diskrete zeitliche Änderung}
\label{sub:diskrete_zeit_aenderung}

Die diskrete zeitliche Ableitung ist:

\[
\delta_t I_k^{(n)} = \frac{I_k^{(n)} - I_k^{(n-1)}}{T},
\]

mit fundamentalem Zeitschrift \( T \).

\section{Diskrete Variation und Euler-Lagrange-Gleichungen}
\label{sec:euler_lagrange_diskret}

Die Dynamik folgt aus dem diskreten Variationsprinzip:

\[
\delta L_I = 0 \quad \text{für alle } \delta I_k^{(n)}.
\]

Die Variation nach \( I_k^{(n)} \) führt zur diskreten Euler-Lagrange-Gleichung:

\[
\delta_t \left( \frac{\partial \mathcal{F}_k}{\partial (\delta_t I_k^{(n)})} \right) + \Delta \left( \frac{\partial \mathcal{F}_k}{\partial (\Delta I_k^{(n)})} \right) - \frac{\partial \mathcal{F}_k}{\partial I_k^{(n)}} = 0.
\]

Hierbei sind:

\begin{itemize}
    \item \( \delta_t f^{(n)} = \frac{f^{(n+1)} - f^{(n)}}{T} \) – diskrete zeitliche Vorwärtsdifferenz,
    \item \( \Delta f_k = \sum_{l \in \mathcal{N}(k)} w_{kl} (f_l - f_k) \) – diskrete räumliche Divergenz.
\end{itemize}

\section{Diskreter Informationsfluss als natürliche Konsequenz}
\label{sec:informationsfluss_konsequenz}

Aus der diskreten Euler-Lagrange-Gleichung folgt unmittelbar der diskrete Informationsfluss:

\[
J_{kl}^{(n)} = \frac{\partial \mathcal{F}_k}{\partial (\Delta_{kl} I^{(n)})},
\]

mit \( \Delta_{kl} I^{(n)} = I_l^{(n)} - I_k^{(n)} \).

Damit wird die diskrete Kontinuitätsgleichung

\[
\delta_t I_k^{(n)} + \sum_{l \in \mathcal{N}(k)} J_{kl}^{(n)} = 0
\]

zu einer direkten Konsequenz des diskreten Variationsprinzips.

\section{Zerlegung in lokale und globale diskrete Beiträge}
\label{sec:lokal_global_zerlegung}

Die Struktur von \( \mathcal{F}_k \) erlaubt eine natürliche Zerlegung:

\[
\mathcal{F}_k = \mathcal{F}_k^{\text{local}} + \mathcal{F}_k^{\text{global}}.
\]

\subsection{Lokaler diskreter Anteil}
\label{sub:lokaler_anteil}

Der lokale Anteil beschreibt lokale Informationsflüsse:

\[
\mathcal{F}_k^{\text{local}} = \alpha \bigl( \delta_t I_k^{(n)} \bigr)^2 + \beta \bigl( \Delta I_k^{(n)} \bigr)^2 + \gamma I_k^{(n)} \ln \left( \frac{I_k^{(n)}}{I_0} \right).
\]

Er führt im diskreten Grenzfall zu:

\begin{itemize}
    \item der diskreten Weber-Kraft,
    \item klassischen Trägheits- und Energiebegriffen,
    \item stabiler numerischer Integration.
\end{itemize}

Die Variation des lokalen Anteils ergibt:

\[
\frac{\partial \mathcal{F}_k^{\text{local}}}{\partial I_k^{(n)}} = \gamma \left( 1 + \ln \left( \frac{I_k^{(n)}}{I_0} \right) \right),
\]
\[
\frac{\partial \mathcal{F}_k^{\text{local}}}{\partial (\delta_t I_k^{(n)})} = 2\alpha \delta_t I_k^{(n)},
\]
\[
\frac{\partial \mathcal{F}_k^{\text{local}}}{\partial (\Delta I_k^{(n)})} = 2\beta \Delta I_k^{(n)}.
\]

\subsection{Globaler diskreter Anteil}
\label{sub:globaler_anteil}

Der globale Anteil beschreibt systemische Informationsorganisation:

\[
\mathcal{F}_k^{\text{global}} = \lambda \frac{(\Delta I_k^{(n)})^2}{I_k^{(n)}} + \mu \frac{\Delta^2 I_k^{(n)}}{I_k^{(n)}},
\]

mit dem diskreten Laplace-Operator:

\[
\Delta^2 I_k^{(n)} = \sum_{l \in \mathcal{N}(k)} w_{kl} \bigl( I_l^{(n)} - I_k^{(n)} \bigr).
\]

Die Variation dieses Terms führt zum diskreten Bohm-Potential:

\[
Q_k^{(n)} = -\frac{\hbar^2}{2m} \frac{\Delta^2 \sqrt{I_k^{(n)}}}{\sqrt{I_k^{(n)}}}.
\]

\section{Diskrete Weber-Kraft und Bohm-Potential als Grenzfälle}
\label{sec:grenzfaelle_weber_bohm}

Die diskrete IWT reproduziert zwei fundamentale Strukturen:

\begin{itemize}
    \item \textbf{Diskrete Weber-Kraft:} Lokaler Grenzfall der diskreten Informationsdynamik,
    \item \textbf{Diskretes Bohm-Potential:} Globaler Grenzfall der diskreten Informationsdynamik.
\end{itemize}

Beide entstehen aus demselben diskreten Funktional – sie sind keine unabhängigen Modelle, sondern zwei Projektionen derselben diskreten Informationsstruktur.

\section{Symmetrien und Erhaltungsgrößen im diskreten Netz}
\label{sec:symmetrien_erhaltung}

Nach dem diskreten Noether-Theorem ergeben sich:

\subsection{Translationsinvarianz}
\label{sub:translationsinvarianz}

Wenn \( \mathcal{F}_k \) nur von Differenzen \( I_k^{(n)} - I_l^{(n)} \) abhängt, dann ist der diskrete Gesamtimpuls erhalten:

\[
P^{(n)} = \sum_k \frac{\partial \mathcal{F}_k}{\partial (\delta_t I_k^{(n)})} = \text{konstant}.
\]

\subsection{Zeitinvarianz}
\label{sub:zeitinvarianz}

Wenn \( \mathcal{F}_k \) nicht explizit von \( n \) abhängt, dann ist die diskrete Energie erhalten:

\[
E^{(n)} = \sum_k \left[ \delta_t I_k^{(n)} \frac{\partial \mathcal{F}_k}{\partial (\delta_t I_k^{(n)})} - \mathcal{F}_k \right] = \text{konstant}.
\]

\subsection{Rotationsinvarianz}
\label{sub:rotationsinvarianz}

Wenn das Netz rotationssymmetrisch ist, dann ist der diskrete Drehimpuls erhalten.

\section{Implementierung als diskreter Update-Algorithmus}
\label{sec:update_algorithmus}

Aus der diskreten Euler-Lagrange-Gleichung ergibt sich eine explizite Update-Regel für \( I_k^{(n+1)} \):

\[
I_k^{(n+1)} = I_k^{(n)} + T \cdot \Phi_k \Bigl( I_k^{(n)}, \{I_l^{(n)}\}, I_k^{(n-1)} \Bigr),
\]

mit

\[
\Phi_k = -\frac{1}{2\alpha} \left[ \Delta \left( 2\beta \Delta I_k^{(n)} \right) - \frac{\partial \mathcal{F}_k}{\partial I_k^{(n)}} \right].
\]

Die konkreten Update-Schritte sind:

\begin{enumerate}
    \item Berechne räumliche Gradienten: \( \Delta I_k^{(n)} \), \( \Delta^2 I_k^{(n)} \),
    \item Berechne partielle Ableitungen: \( \frac{\partial \mathcal{F}_k}{\partial I_k^{(n)}} \), \( \frac{\partial \mathcal{F}_k}{\partial (\Delta I_k^{(n)})} \), etc.,
    \item Löse die Euler-Lagrange-Gleichung nach \( I_k^{(n+1)} \) auf,
    \item Update aller Knoten gleichzeitig (globale Synchronisation).
\end{enumerate}

\section{Grenzfall zur kontinuierlichen Form}
\label{sec:kontinuierlicher_grenzfall_lagrange}

Für \( T \to 0 \) und \( N \to \infty \) mit \( \Delta V_k \to 0 \) ergibt sich der kontinuierliche Grenzfall:

\[
L_I[\rho_I] = \int \mathcal{F}(\rho_I, \nabla \rho_I, \partial_t \rho_I) \, d^3x,
\]

mit \( \rho_I(\vec{r}, t) = \lim I_k^{(n)} \).

Die diskrete Euler-Lagrange-Gleichung geht über in:

\[
\frac{\partial}{\partial t} \left( \frac{\partial \mathcal{F}}{\partial (\partial_t \rho_I)} \right) + \nabla \cdot \left( \frac{\partial \mathcal{F}}{\partial (\nabla \rho_I)} \right) - \frac{\partial \mathcal{F}}{\partial \rho_I} = 0.
\]

\section{Zusammenfassung}
\label{sec:lagrange_zusammenfassung}

Das diskrete Informations-Lagrange-Funktional bildet die mathematische Grundlage der diskreten IWT:

\begin{itemize}
    \item Es beschreibt die Dynamik der diskreten Informationsverteilung \( I_k^{(n)} \).
    \item Es erzeugt die diskrete Kontinuitätsgleichung als natürliche Konsequenz.
    \item Es zerlegt sich in lokale und globale diskrete Beiträge.
    \item Es reproduziert die diskrete Weber-Kraft und das diskrete Bohm-Potential.
    \item Es liefert diskrete Erhaltungssätze aus diskreten Symmetrien.
    \item Es ist algorithmisch direkt implementierbar (siehe Anhang~\ref{app:numerik}).
    \item Der kontinuierliche Grenzfall emergiert für feine Diskretisierung (siehe Anhang~\ref{app:kontinuumslimes}).
\end{itemize}